% !TeX encoding = UTF-8
\documentclass[institutional,screen,none]{ua-engineering-v6-beamer}
% !TeX encoding = UTF-8
\UASetUniversity{Universidad Austral}
\UASetFaculty{Facultad de Ingeniería}
\UASetDepartment{Departamento de Computación}
\UASetCampus{Pilar}
\UASetCareer{Ingeniería Informática}
\UASetCommission{Comisión A}
\UASetProfessor{Equipo docente}
\UASetSemester{Segundo cuatrimestre 2026}
\UASetCourse{Sistemas Distribuidos}
\UASetDate{2026}

\UASetSemester{Segundo cuatrimestre 2026}
\UASetDate{2026}
\usetikzlibrary{decorations.pathreplacing,matrix,backgrounds,shapes.multipart}

\hyphenpenalty=9000
\exhyphenpenalty=9000
\tolerance=1800
\emergencystretch=1.3em

\newcommand{\UASlideTag}[2]{%
  \begin{tikzpicture}[remember picture,overlay]
    \node[anchor=north east,rounded corners=1mm,inner xsep=3pt,inner ysep=1.5pt,
      draw=#2,fill=#2!8,text=#2,font=\fontsize{5.8}{6.4}\selectfont\bfseries]
      at ([xshift=-0.72cm,yshift=-1.08cm]current page.north east) {#1};
  \end{tikzpicture}%
}

\UASetClassNumber{07}
\UASetClassTitle{Logs distribuidos y procesamiento de eventos}
\title{\UACourse}
\subtitle{Clase 07 --- Festival Austral: de dos compras a una historia distribuida y recuperable}
\author{\UAProfessor}
\date{\UADate}

\tikzset{
  uaNode/.style={draw=UAIngNavy,rounded corners=1.5mm,thick,fill=white,
    minimum height=0.70cm,align=center,font=\scriptsize,inner sep=3pt},
  uaSoft/.style={uaNode,fill=UAIngPaper},
  uaOrange/.style={uaNode,draw=UAIngOrange,fill=UAIngPaper},
  uaArrow/.style={->,very thick,>=Latex,draw=UAIngNavy},
  uaOrangeArrow/.style={->,very thick,>=Latex,draw=UAIngOrange},
  uaTiny/.style={font=\tiny,align=center,text=UAIngText},
  uaSmall/.style={font=\scriptsize,align=center,text=UAIngText},
  uaEvent/.style={draw=UAIngNavy,rounded corners=1.2mm,fill=white,
    minimum height=0.44cm,align=center,font=\tiny,inner xsep=3pt,inner ysep=1.5pt},
  uaEventOrange/.style={uaEvent,draw=UAIngOrange,fill=UAIngPaper}
}

\newcommand{\SourceLine}[1]{\vfill{\tiny\color{UAIngMuted}Fuente / referencia: #1}}

\begin{document}
{
\setbeamertemplate{footline}{}
\begin{frame}[plain]
\titlepage
\end{frame}
}

% =============================================================================
% ACTO I: LA NOCHE DEL FESTIVAL Y EL PROBLEMA
% =============================================================================

\begin{frame}{Caso conductor: CampusShop y dos compras simultáneas}
\UASlideTag{NÚCLEO}{UAIngNavy}
\small
\begin{columns}[T]
\column{0.56\textwidth}
\begin{uakeybox}
A las \textbf{11:59} comienza la venta del Festival Austral. Sofía confirma la compra \textbf{9F2A}; casi al mismo tiempo, Tomás inicia la compra \textbf{1C7B}.
\end{uakeybox}
\begin{itemize}
\item Cada compra atraviesa varios estados.
\item Inventario, QR y analítica reaccionan a esos estados.
\item Un componente puede caer y reiniciar.
\item Miles de compras distintas pueden avanzar en paralelo.
\end{itemize}
\column{0.40\textwidth}
\begin{center}
\begin{tikzpicture}[node distance=0.24cm]
\node[uaOrange,text width=3.0cm] (a) {11:59:58\\Sofía: 9F2A};
\node[uaSoft,text width=3.0cm,below=of a] (b) {11:59:59\\Tomás: 1C7B};
\node[uaSoft,text width=3.0cm,below=of b] (c) {12:00:00\\muchos efectos};
\node[uaOrange,text width=3.0cm,below=of c] (d) {12:00:02\\un proceso cae};
\draw[uaOrangeArrow] (a)--(b);
\draw[uaArrow] (b)--(c);
\draw[uaOrangeArrow] (c)--(d);
\end{tikzpicture}
\end{center}
\end{columns}
\end{frame}

\begin{frame}{La promesa visible y el trabajo que queda detrás}
\UASlideTag{NÚCLEO}{UAIngNavy}
\small
\begin{columns}[T]
\column{0.48\textwidth}
\begin{block}{Lo que ve Sofía}
\textbf{Compra confirmada}. Espera que el pedido no desaparezca ni se duplique por un reintento.
\end{block}
\begin{block}{Lo que debe ocurrir después}
Inventario descuenta una unidad, Notificaciones emite el QR y Analítica actualiza el tablero.
\end{block}
\column{0.48\textwidth}
\begin{alertblock}{La tensión}
Si checkout espera sincrónicamente a todos, una dependencia lenta vuelve lenta o fallida a toda la compra.
\end{alertblock}
\begin{uakeybox}
Necesitamos conservar el hecho y permitir que lectores independientes lo procesen a su propio ritmo.
\end{uakeybox}
\end{columns}
\end{frame}

\begin{frame}{De llamadas directas a una historia compartida}
\UASlideTag{NÚCLEO}{UAIngNavy}
\small
\begin{columns}[T]
\column{0.47\textwidth}
\begin{block}{Fan-out sincrónico}
Checkout llama a Inventario, QR y Analítica y espera las respuestas requeridas antes de terminar.
\end{block}
\begin{center}
\begin{tikzpicture}[node distance=0.25cm]
\node[uaOrange,text width=1.55cm] (c) {Checkout};
\node[uaSoft,text width=1.65cm,right=of c,yshift=0.72cm] (i) {Inventario};
\node[uaSoft,text width=1.65cm,right=of c] (q) {QR};
\node[uaSoft,text width=1.65cm,right=of c,yshift=-0.72cm] (a) {Analítica};
\draw[uaArrow] (c)--(i);
\draw[uaArrow] (c)--(q);
\draw[uaArrow] (c)--(a);
\end{tikzpicture}
\end{center}
\column{0.49\textwidth}
\begin{block}{Historia compartida}
Checkout publica \texttt{OrderConfirmed}; cada lector avanza y se recupera de forma independiente.
\end{block}
\begin{center}
\begin{tikzpicture}[node distance=0.25cm]
\node[uaOrange,text width=1.55cm] (c) {Checkout};
\node[uaSoft,text width=1.75cm,right=of c] (l) {Log durable};
\node[uaSoft,text width=1.45cm,right=of l,yshift=0.72cm] (i) {Inventario};
\node[uaSoft,text width=1.45cm,right=of l] (q) {QR};
\node[uaSoft,text width=1.45cm,right=of l,yshift=-0.72cm] (a) {Analítica};
\draw[uaOrangeArrow] (c)--(l);
\draw[uaArrow] (l)--(i);
\draw[uaArrow] (l)--(q);
\draw[uaArrow] (l)--(a);
\end{tikzpicture}
\end{center}
\end{columns}
\end{frame}

\begin{frame}{Base mínima I: quiénes participan}
\UASlideTag{NÚCLEO}{UAIngNavy}
\scriptsize
\begin{center}
\begin{tabular}{p{0.20\textwidth}p{0.64\textwidth}}
\toprule
Actor & Definición de trabajo \\
\midrule
Usuario / cliente & Persona o aplicación que inicia una petición. \\
Balanceador de carga & Distribuye las peticiones entre varias instancias del servicio. \\
Instancia de API & Proceso que valida la petición y ejecuta lógica de negocio. \\
Productor & Código o proceso que publica un evento. \\
Broker & Servidor intermediario que conserva y entrega registros de eventos. \\
Consumidor & Proceso que lee eventos y produce un efecto de negocio. \\
\bottomrule
\end{tabular}
\end{center}
\begin{uakeybox}
\textbf{Base mínima para esta clase:} una petición entra, un productor publica un hecho, uno o varios brokers lo conservan y consumidores independientes lo procesan.
\end{uakeybox}
\end{frame}

\begin{frame}{Base mínima II: proceso y hardware no son lo mismo}
\UASlideTag{NÚCLEO}{UAIngNavy}
\scriptsize
\begin{columns}[T]
\column{0.48\textwidth}
\begin{center}
\begin{tikzpicture}
\node[draw=UAIngNavy,rounded corners=2mm,minimum width=5.15cm,minimum height=3.20cm,fill=UAIngPaper] (n) at (0,0) {};
\node[font=\scriptsize\bfseries,text=UAIngNavy] at (0,1.28) {Nodo físico, VM o host};
\node[uaSoft,text width=1.85cm] (p) at (0,0.48) {Proceso\\software};
\node[uaOrange,text width=1.0cm] (cpu) at (-1.65,-0.70) {CPU};
\node[uaSoft,text width=1.0cm] (ram) at (-0.55,-0.70) {RAM};
\node[uaSoft,text width=1.0cm] (nic) at (0.55,-0.70) {NIC};
\node[uaOrange,text width=1.0cm] (ssd) at (1.65,-0.70) {SSD};
\draw[uaArrow] (p)--(cpu);
\draw[uaArrow] (p)--(ram);
\draw[uaArrow] (p)--(nic);
\draw[uaArrow] (p)--(ssd);
\end{tikzpicture}
\end{center}
\column{0.48\textwidth}
\begin{uainfo}
\textbf{CPU:} serializa, comprime y ejecuta lógica.\\[0.07cm]
\textbf{RAM:} buffers, lotes y estado temporal.\\[0.07cm]
\textbf{NIC:} interfaz de red.\\[0.07cm]
\textbf{SSD:} segmentos de log y estado durable.
\end{uainfo}
\begin{alertblock}{Distinción}
La API, el broker y el consumidor son procesos; CPU, RAM, NIC y SSD son recursos del nodo.
\end{alertblock}
\end{columns}
\end{frame}

\begin{frame}{Base mínima III: registro, evento y log}
\UASlideTag{NÚCLEO}{UAIngNavy}
\scriptsize
\begin{center}
\begin{tabular}{p{0.20\textwidth}p{0.64\textwidth}}
\toprule
Término & Definición de trabajo \\
\midrule
Registro & Unidad guardada en una secuencia: posee identidad, datos y posición. \\
Evento & Registro que describe un hecho ya ocurrido, como \texttt{OrderConfirmed}. \\
Append & Agregar un registro al final sin reescribir toda la historia. \\
Log & Secuencia retenida y ordenada que puede releerse. \\
Replay & Volver a leer registros anteriores para reconstruir estado o procesamiento. \\
\bottomrule
\end{tabular}
\end{center}
\begin{uakeybox}
Aquí “log” no significa solo texto de depuración: es una estructura de datos que conserva historia.
\end{uakeybox}
\end{frame}

\begin{frame}{Vocabulario previo: topic, partición, key y offset}
\UASlideTag{NÚCLEO}{UAIngNavy}
\scriptsize
\begin{center}
\begin{tabular}{p{0.19\textwidth}p{0.65\textwidth}}
\toprule
Término & Significado \\
\midrule
Topic & Canal lógico que reúne eventos de un mismo contrato, por ejemplo \texttt{orders.lifecycle}. \\
Partición & Subsecuencia ordenada del topic; se almacena, replica y consume como unidad. \\
Key & Valor estable que determina a qué partición va un evento. \\
Offset & Posición creciente de un registro dentro de una partición. \\
\bottomrule
\end{tabular}
\end{center}
\begin{alertblock}{Dos límites que no deben confundirse}
El orden existe \textbf{dentro de una partición}; el offset identifica posición, no identidad de negocio.
\end{alertblock}
\end{frame}

\begin{frame}{El accidente de las 12:00:02}
\UASlideTag{PREDICCIÓN}{UAIngOrange}
\small
\begin{center}
\begin{tikzpicture}[node distance=0.32cm]
\node[uaOrange,text width=2.25cm] (co) {Checkout\\confirmó 9F2A};
\node[uaSoft,text width=2.20cm,right=of co,yshift=0.85cm] (i) {Inventario\\procesa};
\node[uaSoft,text width=2.20cm,right=of co] (q) {QR\\cae 40 s};
\node[uaSoft,text width=2.20cm,right=of co,yshift=-0.85cm] (a) {Analítica\\se atrasa};
\draw[uaArrow] (co)--(i);
\draw[uaOrangeArrow] (co)--(q);
\draw[uaArrow] (co)--(a);
\end{tikzpicture}
\end{center}
\begin{alertblock}{Predicción}
¿Cómo recupera QR el trabajo? ¿Qué debe conservarse para que no dependa de que checkout siga vivo?
\end{alertblock}
\end{frame}

\begin{frame}{Pregunta central: del hecho al efecto recuperable}
\UASlideTag{NÚCLEO}{UAIngNavy}
\small
\begin{uakeybox}
¿Cómo convertir dos compras concurrentes en una historia durable que varios componentes puedan procesar y reprocesar sin perder progreso ni repetir el efecto lógico?
\end{uakeybox}
\begin{center}
\begin{tikzpicture}[node distance=0.28cm]
\node[uaOrange,text width=2.1cm] (h) {Hechos\\9F2A y 1C7B};
\node[uaSoft,text width=2.1cm,right=of h] (s) {Historias\\ordenadas};
\node[uaSoft,text width=2.1cm,right=of s] (r) {Lectores\\independientes};
\node[uaOrange,text width=2.1cm,right=of r] (e) {Efectos\\correctos};
\draw[uaOrangeArrow] (h)--(s);
\draw[uaArrow] (s)--(r);
\draw[uaOrangeArrow] (r)--(e);
\end{tikzpicture}
\end{center}
\end{frame}

\begin{frame}{Evidencias de aprendizaje de esta clase}
\UASlideTag{NÚCLEO}{UAIngNavy}
\scriptsize
\begin{columns}[T]
\column{0.48\textwidth}
\begin{block}{Durante la clase}
\begin{itemize}
\item completar las historias 9F2A y 1C7B;
\item elegir y justificar una key;
\item ubicar ambas historias en tres particiones;
\item mover un crash en el timeline de consumo;
\item diseñar outbox e inbox.
\end{itemize}
\end{block}
\column{0.48\textwidth}
\begin{block}{Al cerrar}
\begin{itemize}
\item declarar qué orden se preserva y cuál no;
\item explicar lag y rebalanceo desde la historia;
\item delimitar una garantía exactly-once;
\item relacionar síntoma, nodo y recurso físico.
\end{itemize}
\end{block}
\end{columns}
\end{frame}

\begin{frame}{La historia de hoy en cinco capítulos}
\UASlideTag{NÚCLEO}{UAIngNavy}
\scriptsize
\begin{center}
\begin{tikzpicture}[node distance=0.20cm]
\node[uaOrange,text width=1.65cm] (a) {1. Publicar\\el hecho};
\node[uaSoft,text width=1.65cm,right=of a] (b) {2. Ordenar\\cada historia};
\node[uaSoft,text width=1.65cm,right=of b] (c) {3. Consumir\\y recuperar};
\node[uaSoft,text width=1.65cm,right=of c] (d) {4. Publicar\\sin dual write};
\node[uaOrange,text width=1.65cm,right=of d] (e) {5. Operar\\el stream};
\draw[uaOrangeArrow] (a)--(b);
\draw[uaArrow] (b)--(c);
\draw[uaArrow] (c)--(d);
\draw[uaOrangeArrow] (d)--(e);
\end{tikzpicture}
\end{center}
\begin{uakeybox}
Los conceptos aparecerán cuando la historia los necesite; después volveremos a 9F2A y 1C7B para verificar que el recorrido quedó completo.
\end{uakeybox}
\end{frame}

\begin{frame}{Diagnóstico: razone antes de elegir tecnología}
\UASlideTag{ACTIVIDAD}{UAIngOrange}
\small
\begin{enumerate}
\item ¿Los eventos de una misma compra deben conservar un orden?
\item ¿Los eventos de compras distintas necesitan un orden global único?
\item Si Inventario cae después de descontar stock, ¿qué debe reaparecer al reiniciar?
\item Si el checkout guardó la compra pero no publicó el hecho, ¿cómo se entera QR?
\item Si una partición acumula trabajo, ¿agregar consumidores siempre la divide?
\end{enumerate}
\begin{alertblock}{Consigna}
Responda con el comportamiento deseado; todavía no use Kafka, offset, outbox o rebalanceo.
\end{alertblock}
\end{frame}

\begin{frame}{Puente desde la clase 6: del WAL local al log compartido}
\UASlideTag{NÚCLEO}{UAIngNavy}
\small
\begin{columns}[T]
\column{0.47\textwidth}
\begin{block}{WAL local}
Historia interna de un motor para recovery de sus páginas y transacciones.
\end{block}
\column{0.47\textwidth}
\begin{block}{Log distribuido}
Historia compartida y retenida que varios procesos leen con progreso independiente.
\end{block}
\end{columns}
\begin{uakeybox}
La idea de append y replay continúa; cambian la audiencia, la escala, la distribución y la semántica del efecto.
\end{uakeybox}
\end{frame}

\begin{frame}{Actividad: ¿qué debe ocurrir ahora y qué puede ocurrir después?}
\UASlideTag{ACTIVIDAD}{UAIngOrange}
\small
Clasifique cada efecto del pedido 9F2A:
\begin{center}
\begin{tabular}{p{0.33\textwidth}p{0.23\textwidth}p{0.27\textwidth}}
\toprule
Efecto & ¿Antes del OK? & Razón \\
\midrule
Confirmar pedido & & \\
Reservar/descontar stock & & \\
Enviar QR & & \\
Actualizar analítica & & \\
Recalcular ranking & & \\
\bottomrule
\end{tabular}
\end{center}
\begin{alertblock}{Objetivo}
Separar el camino sincrónico obligatorio de los efectos que pueden ocurrir asincrónicamente.
\end{alertblock}
\end{frame}

% =============================================================================
\UASectionSlide{Capítulo 1 --- publicar una historia compartida}
% =============================================================================

\begin{frame}{Arquitectura mínima: productor, brokers y consumidores}
\UASlideTag{NÚCLEO}{UAIngNavy}
\small
\begin{center}
\begin{tikzpicture}[node distance=0.34cm]
\node[uaOrange,text width=2.15cm] (p) {Checkout\\productor};
\node[uaSoft,text width=2.35cm,right=of p] (b) {Cluster de brokers\\log retenido};
\node[uaSoft,text width=1.85cm,right=of b,yshift=0.82cm] (i) {Inventario};
\node[uaSoft,text width=1.85cm,right=of b] (q) {QR};
\node[uaSoft,text width=1.85cm,right=of b,yshift=-0.82cm] (a) {Analítica};
\draw[uaOrangeArrow] (p)--node[above,uaTiny]{publica}(b);
\draw[uaArrow] (b)--(i);
\draw[uaArrow] (b)--(q);
\draw[uaArrow] (b)--(a);
\end{tikzpicture}
\end{center}
\begin{uakeybox}
Una publicación agrega un hecho a la historia; no obliga a que todos los consumidores terminen al mismo tiempo.
\end{uakeybox}
\end{frame}

\begin{frame}{Kafka: una implementación real del modelo}
\UASlideTag{CASO REAL}{UAIngOrange}
\small
\begin{uainfo}
\textbf{Apache Kafka} es una plataforma distribuida de eventos. Sus brokers conservan logs particionados; productores publican y consumidores leen con progreso propio.
\end{uainfo}
\begin{itemize}
\item Topic: categoría lógica de eventos.
\item Partición: historia ordenada y escalable.
\item Broker: proceso que aloja particiones.
\item Consumer group: instancias que se reparten particiones.
\end{itemize}
\begin{alertblock}{Método de la clase}
Primero razonamos con el modelo; luego usamos Kafka como implementación concreta, no como sustituto de los conceptos.
\end{alertblock}
\end{frame}

\begin{frame}{Del software al hardware: recorrido de un evento}
\UASlideTag{EXTENSIÓN}{UAIngNavy}
\scriptsize
\begin{center}
\begin{tabular}{p{0.17\textwidth}p{0.25\textwidth}p{0.36\textwidth}}
\toprule
Nodo & Proceso & Recursos involucrados \\
\midrule
API & productor & CPU serializa; RAM forma lote; NIC envía \\
Broker líder & broker & NIC recibe; CPU valida; RAM/page cache; SSD conserva \\
Broker seguidor & réplica & NIC recibe copia; CPU verifica; SSD conserva \\
Consumidor & lector & NIC trae lote; CPU procesa; RAM/DB guarda estado \\
\bottomrule
\end{tabular}
\end{center}
\begin{uakeybox}
La abstracción “publicar un evento” se materializa como cómputo, buffers, paquetes y escrituras en varios nodos.
\end{uakeybox}
\end{frame}

% =============================================================================
\UASectionSlide{Capítulo 2 --- dos historias, tres particiones}
% =============================================================================

\begin{frame}{Primero complete las dos historias}
\UASlideTag{NÚCLEO}{UAIngNavy}
\small
\begin{columns}[T]
\column{0.48\textwidth}
\begin{block}{Compra 9F2A — Sofía}
\begin{enumerate}
\item \texttt{OrderCreated}
\item \texttt{PaymentAuthorized}
\item \texttt{OrderConfirmed}
\item \texttt{QrIssued}
\end{enumerate}
\end{block}
\column{0.48\textwidth}
\begin{block}{Compra 1C7B — Tomás}
\begin{enumerate}
\item \texttt{OrderCreated}
\item \texttt{PaymentRejected}
\item \texttt{OrderCancelled}
\end{enumerate}
\end{block}
\end{columns}
\begin{alertblock}{Pregunta}
¿Qué pares de eventos exigen orden relativo y qué eventos pueden intercalarse libremente?
\end{alertblock}
\end{frame}

\begin{frame}{Causalidad necesaria: qué orden no podemos perder}
\UASlideTag{NÚCLEO}{UAIngNavy}
\small
\begin{uainfo}
Dos eventos tienen una dependencia causal cuando el segundo presupone que el primero ya ocurrió.
\end{uainfo}
\begin{center}
\begin{tikzpicture}[node distance=0.24cm]
\node[uaOrange,text width=1.55cm] (a) {Created\\9F2A};
\node[uaSoft,text width=1.75cm,right=of a] (b) {Payment\\Authorized};
\node[uaSoft,text width=1.65cm,right=of b] (c) {Order\\Confirmed};
\node[uaOrange,text width=1.45cm,right=of c] (d) {QR\\Issued};
\draw[uaOrangeArrow] (a)--(b);
\draw[uaArrow] (b)--(c);
\draw[uaOrangeArrow] (c)--(d);
\end{tikzpicture}
\end{center}
\begin{itemize}
\item Para 9F2A, QR no debe aparecer antes de la confirmación.
\item Para 1C7B, cancelación no debe aparecer antes del rechazo del pago.
\item Entre 9F2A y 1C7B no necesitamos un orden global único.
\end{itemize}
\end{frame}

\begin{frame}{Qué es una key de particionamiento}
\UASlideTag{NÚCLEO}{UAIngNavy}
\small
\begin{uakeybox}
La \textbf{key} es un valor estable extraído del evento. El productor aplica una función determinista para elegir la partición.
\end{uakeybox}
\[
\text{partición}=h(\text{key})\bmod P
\]
\begin{center}
\scriptsize $h$: función hash determinista; $P$: cantidad de particiones.
\end{center}
\begin{columns}[T]
\column{0.48\textwidth}
\begin{block}{Misma key}
Todos los eventos de esa entidad llegan a la misma partición y pueden conservar su orden local.
\end{block}
\column{0.48\textwidth}
\begin{block}{Keys distintas}
Pueden llegar a particiones diferentes y procesarse en paralelo.
\end{block}
\end{columns}
\end{frame}

\begin{frame}{Elección para la historia: key = order\_id}
\UASlideTag{NÚCLEO}{UAIngNavy}
\small
\begin{center}
\begin{tikzpicture}[node distance=0.45cm]
\node[uaOrange,text width=2.25cm] (a) {Eventos 9F2A\\key=9F2A};
\node[uaSoft,text width=2.15cm,right=of a] (ha) {$h(9F2A)\bmod 3$};
\node[uaOrange,text width=1.55cm,right=of ha] (p2) {P2};
\node[uaSoft,text width=2.25cm,below=0.70cm of a] (b) {Eventos 1C7B\\key=1C7B};
\node[uaSoft,text width=2.15cm,right=of b] (hb) {$h(1C7B)\bmod 3$};
\node[uaSoft,text width=1.55cm,right=of hb] (p0) {P0};
\draw[uaOrangeArrow] (a)--(ha);
\draw[uaOrangeArrow] (ha)--(p2);
\draw[uaArrow] (b)--(hb);
\draw[uaArrow] (hb)--(p0);
\end{tikzpicture}
\end{center}
\begin{uakeybox}
No elegimos manualmente la partición para cada evento: elegimos una key estable y la función produce siempre el mismo destino. Los valores P0 y P2 son ilustrativos para esta historia.
\end{uakeybox}
\end{frame}

\begin{frame}{Las dos historias ya ubicadas en tres particiones}
\UASlideTag{NÚCLEO}{UAIngNavy}
\scriptsize
\begin{center}
\begin{tabular}{p{0.12\textwidth}p{0.72\textwidth}}
\toprule
Partición & Registros en orden de offset \\
\midrule
P0 & 105:Created(1C7B) $\rightarrow$ 106:PaymentRejected(1C7B) $\rightarrow$ 107:Cancelled(1C7B) \\
P1 & 221:Created(7D8E) $\rightarrow$ 222:Confirmed(7D8E) $\rightarrow$ otros pedidos \\
P2 & 417:Created(9F2A) $\rightarrow$ 418:PaymentAuthorized(9F2A) $\rightarrow$ 419:Confirmed(9F2A) $\rightarrow$ 420:QrIssued(9F2A) \\
\bottomrule
\end{tabular}
\end{center}
\begin{alertblock}{Resultado}
Se conserva el orden causal de cada pedido. No existe un offset global que ordene P0 frente a P2.
\end{alertblock}
\end{frame}

\begin{frame}{Qué ganamos y a qué renunciamos}
\UASlideTag{NÚCLEO}{UAIngNavy}
\small
\begin{columns}[T]
\column{0.47\textwidth}
\begin{block}{Ganamos}
\begin{itemize}
\item orden por pedido;
\item paralelismo entre pedidos;
\item distribución de escritura y consumo;
\item recuperación por partición.
\end{itemize}
\end{block}
\column{0.49\textwidth}
\begin{block}{Renunciamos}
\begin{itemize}
\item orden global entre todas las compras;
\item procesamiento paralelo dentro de una misma key;
\item rebalancear una key sin mover su historia.
\end{itemize}
\end{block}
\end{columns}
\begin{uakeybox}
La partición es simultáneamente una frontera de orden, capacidad y paralelismo.
\end{uakeybox}
\end{frame}

\begin{frame}{Actividad: reconstruya el reparto y defienda la key}
\UASlideTag{ACTIVIDAD}{UAIngOrange}
\small
\begin{enumerate}
\item Coloque los siete eventos de 9F2A y 1C7B en P0, P1 o P2.
\item Mantenga juntos los eventos de cada pedido.
\item Asigne offsets crecientes en cada partición.
\item Dibuje dos intercalados globales distintos que sigan siendo correctos.
\item Explique qué cambiaría si la key fuera \texttt{event\_id}.
\end{enumerate}
\begin{alertblock}{Evidencia esperada}
La respuesta debe declarar el orden preservado, el paralelismo obtenido y el orden global que deja de existir.
\end{alertblock}
\end{frame}

\begin{frame}{Qué hace buena a una key}
\UASlideTag{NÚCLEO}{UAIngNavy}
\scriptsize
\begin{center}
\begin{tabular}{p{0.23\textwidth}p{0.60\textwidth}}
\toprule
Criterio & Pregunta de diseño \\
\midrule
Causalidad & ¿Agrupa exactamente los eventos cuyo orden relativo importa? \\
Estabilidad & ¿El valor permanece igual durante toda la vida de la historia? \\
Cardinalidad & ¿Existen suficientes valores distintos para repartir la carga? \\
Granularidad & ¿Evita serializar eventos que no necesitan esperar entre sí? \\
Distribución & ¿Ningún valor concentra una fracción desproporcionada del tráfico? \\
Observabilidad & ¿Podemos medir bytes, QPS y lag por key/partición? \\
\bottomrule
\end{tabular}
\end{center}
\begin{uakeybox}
Una buena key es el identificador estable más pequeño que preserva la causalidad necesaria y distribuye historias independientes.
\end{uakeybox}
\end{frame}

\begin{frame}{Tres keys posibles para el mismo evento}
\UASlideTag{NÚCLEO}{UAIngNavy}
\scriptsize
\begin{center}
\begin{tabular}{p{0.20\textwidth}p{0.30\textwidth}p{0.32\textwidth}}
\toprule
Key & Efecto sobre el orden & Efecto sobre la carga \\
\midrule
\texttt{order\_id} & conserva el ciclo de vida del pedido & muchos pedidos se reparten \\
\texttt{event\_id} & eventos del mismo pedido pueden separarse & excelente reparto, causalidad rota \\
\texttt{campus\_id} & ordena todo el campus & una venta masiva concentra todo \\
\bottomrule
\end{tabular}
\end{center}
\begin{alertblock}{No existe una key universal}
La key depende del invariante y del orden que necesita ese topic. Otra vista puede requerir otro topic o una proyección derivada.
\end{alertblock}
\end{frame}

\begin{frame}{Hot key y hot partition: el límite de una historia muy popular}
\UASlideTag{EXTENSIÓN}{UAIngNavy}
\small
\begin{uainfo}
Una \textbf{hot key} produce mucho más tráfico que las demás. Como una key pertenece a una sola partición, puede convertirla en \textbf{hot partition}.
\end{uainfo}
\begin{itemize}
\item Un solo pedido normal genera pocos eventos: \texttt{order\_id} funciona bien.
\item Una key global como \texttt{campus\_id} concentra toda la venta.
\item Una historia individual extremadamente activa no puede paralelizarse sin redefinir su frontera de orden.
\end{itemize}
\begin{alertblock}{Opciones}
Rediseñar la key, dividir la entidad si la semántica lo permite, crear agregación jerárquica o aceptar el cuello de botella para preservar orden.
\end{alertblock}
\end{frame}

\begin{frame}{De particiones lógicas a brokers físicos}
\UASlideTag{NÚCLEO}{UAIngNavy}
\small
\begin{center}
\begin{tikzpicture}[node distance=0.50cm]
\node[uaOrange,text width=2.25cm] (a) {Broker A\\líder P0};
\node[uaSoft,text width=2.25cm,right=of a] (b) {Broker B\\líder P1};
\node[uaSoft,text width=2.25cm,right=of b] (c) {Broker C\\líder P2};
\node[uaTiny,below=0.20cm of a] {copias P1, P2};
\node[uaTiny,below=0.20cm of b] {copias P0, P2};
\node[uaTiny,below=0.20cm of c] {copias P0, P1};
\end{tikzpicture}
\end{center}
\begin{itemize}
\item P0 contiene la historia 1C7B; su líder atiende appends y lecturas.
\item P2 contiene la historia 9F2A; su líder está en otro nodo.
\item Las copias protegen frente a pérdida de nodo, pero consumen red y disco.
\end{itemize}
\end{frame}

\begin{frame}{Publicar Confirmed(9F2A): del productor al ACK}
\UASlideTag{NÚCLEO}{UAIngNavy}
\scriptsize
\begin{center}
\begin{tikzpicture}[node distance=0.55cm and 0.85cm]
\node[uaOrange,text width=1.85cm] (p) at (-3.7,0.75) {Productor\\batch en RAM};
\node[uaSoft,text width=2.00cm] (leader) at (-0.75,0.75) {Broker C\\líder de P2};
\node[uaOrange,text width=1.35cm] (ack) at (2.25,0.75) {ACK};
\node[uaSoft,text width=1.95cm] (r1) at (-1.95,-0.75) {Broker A\\copia de P2};
\node[uaSoft,text width=1.95cm] (r2) at (0.45,-0.75) {Broker B\\copia de P2};
\draw[uaOrangeArrow] (p)--node[above,uaTiny]{append}(leader);
\draw[uaArrow] (leader)--node[above,uaTiny]{política cumplida}(ack);
\draw[uaArrow] (leader.south west)--node[left,uaTiny]{réplica}(r1.north);
\draw[uaArrow] (leader.south east)--node[right,uaTiny]{réplica}(r2.north);
\end{tikzpicture}
\end{center}
\begin{uakeybox}
El líder responde cuando se cumple la política configurada de confirmación. El ACK no significa automáticamente “un \texttt{fsync} independiente por evento”.
\end{uakeybox}
\end{frame}

\begin{frame}{Timeout del productor: el evento 419 puede o no existir}
\UASlideTag{NÚCLEO}{UAIngNavy}
\small
\begin{columns}[T]
\column{0.48\textwidth}
\begin{block}{Historia A}
\texttt{Confirmed(9F2A)} nunca llegó a P2. Reintentar es necesario.
\end{block}
\column{0.48\textwidth}
\begin{block}{Historia B}
P2 agregó el evento en offset 419, pero se perdió el ACK. Reintentar puede duplicarlo.
\end{block}
\end{columns}
\begin{uakeybox}
La idempotencia del productor usa identidad y secuencia para que el broker reconozca reenvíos del mismo batch. No deduplica el cobro ni el email externo.
\end{uakeybox}
\end{frame}

% =============================================================================
\UASectionSlide{Capítulo 3 --- varios lectores y progreso recuperable}
% =============================================================================

\begin{frame}{Consumidor, poll y committed offset}
\UASlideTag{NÚCLEO}{UAIngNavy}
\scriptsize
\begin{center}
\begin{tabular}{p{0.21\textwidth}p{0.63\textwidth}}
\toprule
Término & Significado en la historia \\
\midrule
\texttt{poll} & llamada con la que el consumidor solicita un lote al broker \\
Polling & repetición de esa consulta para seguir leyendo \\
Position & próximo offset que la instancia viva intentará leer \\
committed offset & progreso durable desde el cual otra instancia reinicia \\
\bottomrule
\end{tabular}
\end{center}
\begin{uakeybox}
El broker entrega registros; Inventario decide cómo transformar cada evento en stock.
\end{uakeybox}
\end{frame}

\begin{frame}{Consumer group: repartir las tres particiones}
\UASlideTag{NÚCLEO}{UAIngNavy}
\small
\begin{center}
\begin{tikzpicture}[node distance=0.45cm]
\node[uaOrange,text width=1.40cm] (p0) {P0\\1C7B};
\node[uaSoft,text width=1.40cm,right=of p0] (p1) {P1\\otros};
\node[uaSoft,text width=1.40cm,right=of p1] (p2) {P2\\9F2A};
\node[uaSoft,text width=1.70cm,below=1.0cm of p0] (c1) {Inventario-1};
\node[uaSoft,text width=1.70cm,below=1.0cm of p1] (c2) {Inventario-2};
\node[uaSoft,text width=1.70cm,below=1.0cm of p2] (c3) {Inventario-3};
\draw[uaArrow] (p0)--(c1);
\draw[uaArrow] (p1)--(c2);
\draw[uaArrow] (p2)--(c3);
\end{tikzpicture}
\end{center}
\begin{uakeybox}
Dentro de un grupo, una partición se asigna a una sola instancia a la vez. El paralelismo activo no supera el número de particiones.
\end{uakeybox}
\end{frame}

\begin{frame}{Las dos historias vistas por Inventario}
\UASlideTag{NÚCLEO}{UAIngNavy}
\scriptsize
\begin{columns}[T]
\column{0.48\textwidth}
\begin{block}{Inventario-1 lee P0}
\texttt{Created(1C7B)}\\
\texttt{PaymentRejected(1C7B)}\\
\texttt{Cancelled(1C7B)}
\end{block}
\column{0.48\textwidth}
\begin{block}{Inventario-3 lee P2}
\texttt{Created(9F2A)}\\
\texttt{PaymentAuthorized(9F2A)}\\
\texttt{Confirmed(9F2A)}\\
\texttt{QrIssued(9F2A)}
\end{block}
\end{columns}
\begin{alertblock}{Consecuencia}
Los dos pedidos se procesan en paralelo; dentro de cada pedido se conserva la secuencia elegida por la key.
\end{alertblock}
\end{frame}

\begin{frame}{El accidente de Inventario-3}
\UASlideTag{PREDICCIÓN}{UAIngOrange}
\small
Inventario-3 lee \texttt{Confirmed(9F2A)} en P2:419, descuenta stock y cae antes de guardar su committed offset.
\begin{center}
\begin{tikzpicture}[node distance=0.34cm]
\node[uaOrange,text width=1.55cm] (r) {poll\\419};
\node[uaSoft,text width=1.65cm,right=of r] (e) {stock\\--1};
\node[uaOrange,text width=1.45cm,right=of e] (x) {CRASH};
\node[uaSoft,text width=1.80cm,right=of x] (s) {restart\\desde 419};
\draw[uaOrangeArrow] (r)--(e);
\draw[uaArrow] (e)--(x);
\draw[uaOrangeArrow] (x)--(s);
\end{tikzpicture}
\end{center}
\begin{alertblock}{Predicción}
¿Qué reaparece y cómo evitamos descontar stock por segunda vez?
\end{alertblock}
\end{frame}

\begin{frame}{At-most-once: progreso antes que efecto}
\UASlideTag{NÚCLEO}{UAIngNavy}
\small
\[
\text{poll 419}\rightarrow\text{commit 420}\rightarrow\text{crash}\rightarrow\text{efecto nunca aplicado}
\]
\begin{block}{Resultado}
El evento no reaparece, pero se puede perder el descuento de stock.
\end{block}
\begin{uakeybox}
At-most-once prioriza no repetir; una falla entre progreso y efecto puede producir pérdida lógica.
\end{uakeybox}
\end{frame}

\begin{frame}{At-least-once: efecto antes que progreso}
\UASlideTag{NÚCLEO}{UAIngNavy}
\small
\[
\text{poll 419}\rightarrow\text{stock --1}\rightarrow\text{crash}\rightarrow\text{replay 419}
\]
\begin{block}{Resultado}
El evento reaparece; si el efecto no es idempotente, el stock puede descontarse dos veces.
\end{block}
\begin{uakeybox}
At-least-once facilita progreso frente a fallas, pero exige idempotencia o deduplicación del efecto.
\end{uakeybox}
\end{frame}

\begin{frame}{Inbox idempotente: recordar el evento junto con el efecto}
\UASlideTag{NÚCLEO}{UAIngNavy}
\small
\begin{center}
\begin{tikzpicture}[node distance=0.30cm]
\node[uaOrange,text width=1.70cm] (ev) {event\_id\\evt-9F2A-3};
\node[uaSoft,text width=2.10cm,right=of ev] (tx) {Transacción local\\inbox + stock};
\node[uaSoft,text width=1.75cm,right=of tx] (db) {DB durable};
\node[uaOrange,text width=1.55cm,right=of db] (off) {commit\\offset};
\draw[uaOrangeArrow] (ev)--(tx);
\draw[uaArrow] (tx)--(db);
\draw[uaOrangeArrow] (db)--(off);
\end{tikzpicture}
\end{center}
\begin{uakeybox}
Si 419 reaparece, la inbox demuestra que el efecto ya fue aplicado. Identidad y efecto deben confirmarse atómicamente.
\end{uakeybox}
\end{frame}

\begin{frame}{Actividad: mueva el crash en la historia 9F2A}
\UASlideTag{ACTIVIDAD}{UAIngOrange}
\small
Ubique el crash en cuatro cortes:
\begin{enumerate}
\item antes de aplicar stock;
\item después de stock, antes de inbox;
\item después de inbox, antes de committed offset;
\item después del committed offset.
\end{enumerate}
Para cada corte, responda:
\begin{itemize}
\item ¿qué estado es durable?;
\item ¿qué offset reaparece?;
\item ¿se pierde o duplica el efecto?;
\item ¿qué debe ser atómico?
\end{itemize}
\end{frame}

\begin{frame}{Rebalanceo: Inventario-3 cae y P2 cambia de dueño}
\UASlideTag{NÚCLEO}{UAIngNavy}
\small
\begin{center}
\begin{tikzpicture}[node distance=0.55cm]
\node[uaOrange,text width=1.50cm] (p2) {P2\\9F2A};
\node[uaSoft,text width=1.85cm,right=of p2] (c3) {Inventario-3\\caído};
\node[uaSoft,text width=1.85cm,below=0.90cm of c3] (c2) {Inventario-2\\nuevo dueño};
\draw[uaOrangeArrow] (p2)--node[above,uaTiny]{antes}(c3);
\draw[uaArrow] (p2) to[bend right=18] node[below,uaTiny]{reasignación}(c2);
\end{tikzpicture}
\end{center}
\begin{uainfo}
\textbf{Rebalanceo:} proceso por el cual el grupo reasigna particiones cuando cambia la membresía o la topología.
\end{uainfo}
\begin{itemize}
\item Inventario-2 reinicia P2 desde el committed offset.
\item Durante la reasignación puede haber pausa y crecimiento temporal del lag.
\end{itemize}
\end{frame}

\begin{frame}{Lag: P2 avanza más rápido que Inventario-2}
\UASlideTag{NÚCLEO}{UAIngNavy}
\small
\[
L_o=O_{\text{final de P2}}-O_{\text{confirmado por Inventario}}
\]
\begin{uainfo}
\textbf{Lag en offsets:} cantidad de registros pendientes.\\
\textbf{Lag temporal:} edad del evento pendiente más antiguo.
\end{uainfo}
\begin{alertblock}{Hilo narrativo}
El lag crece porque P2 sigue recibiendo pedidos mientras Inventario-2 restaura estado, consulta una base lenta o procesa menos eventos por segundo.
\end{alertblock}
\end{frame}

\begin{frame}{Dashboard por partición: no diagnostique solo el promedio}
\UASlideTag{EXTENSIÓN}{UAIngNavy}
\scriptsize
\begin{center}
\begin{tabular}{p{0.13\textwidth}rrrrp{0.22\textwidth}}
\toprule
Partición & Ingreso/s & Proceso/s & Lag & CPU & Hipótesis \\
\midrule
P0 & 220 & 240 & 0 & 45\% & sana \\
P1 & 210 & 230 & 12 & 48\% & fluctuación \\
P2 & 520 & 260 & 19.400 & 91\% & carga/skew/restore \\
\bottomrule
\end{tabular}
\end{center}
\begin{uakeybox}
Una hot partition, un restore o una dependencia downstream lenta pueden producir el mismo síntoma superficial: lag creciente.
\end{uakeybox}
\end{frame}

% =============================================================================
\UASectionSlide{Capítulo 4 --- publicar el hecho sin dual write}
% =============================================================================

\begin{frame}{La nueva falla: pedido confirmado, evento ausente}
\UASlideTag{NÚCLEO}{UAIngNavy}
\small
Checkout confirma 9F2A en la base y luego intenta publicar \texttt{OrderConfirmed}.
\begin{center}
\begin{tikzpicture}[node distance=0.35cm]
\node[uaOrange,text width=1.80cm] (api) {Checkout};
\node[uaSoft,text width=2.00cm,right=of api] (db) {DB Orders\\COMMIT};
\node[uaOrange,text width=1.55cm,right=of db] (x) {CRASH};
\node[uaSoft,text width=2.00cm,right=of x] (br) {Broker\\sin evento};
\draw[uaOrangeArrow] (api)--(db);
\draw[uaArrow] (db)--(x);
\draw[uaArrow,dashed] (x)--(br);
\end{tikzpicture}
\end{center}
\begin{alertblock}{Consecuencia}
La compra existe, pero QR, Analítica e Inventario nunca reciben el hecho.
\end{alertblock}
\end{frame}

\begin{frame}{Dual write: ningún orden elimina todas las ventanas}
\UASlideTag{NÚCLEO}{UAIngNavy}
\small
\begin{columns}[T]
\column{0.48\textwidth}
\begin{block}{Base primero}
DB confirma $\rightarrow$ crash $\rightarrow$ evento ausente.
\end{block}
\column{0.48\textwidth}
\begin{block}{Broker primero}
Evento publicado $\rightarrow$ crash/rollback $\rightarrow$ consumidores ven una compra inexistente.
\end{block}
\end{columns}
\begin{uakeybox}
Dos escrituras independientes no forman una transacción atómica solo por elegir un orden.
\end{uakeybox}
\end{frame}

\begin{frame}{Transactional outbox: negocio e intención en una transacción}
\UASlideTag{NÚCLEO}{UAIngNavy}
\small
\begin{uainfo}
Una \textbf{outbox transaccional} es una tabla o colección de eventos pendientes escrita en la misma transacción local que el cambio de negocio.
\end{uainfo}
\begin{center}
\begin{tikzpicture}[node distance=0.30cm]
\node[uaOrange,text width=1.75cm] (api) {Checkout};
\node[uaSoft,text width=2.30cm,right=of api] (tx) {Transacción local\\pedido + outbox};
\node[uaSoft,text width=1.80cm,right=of tx] (db) {DB durable};
\draw[uaOrangeArrow] (api)--(tx);
\draw[uaArrow] (tx)--(db);
\end{tikzpicture}
\end{center}
\begin{uakeybox}
Si el COMMIT local existe, también existe una intención durable de publicar \texttt{OrderConfirmed(9F2A)}.
\end{uakeybox}
\end{frame}

\begin{frame}{Relay, polling y CDC: cómo sale el evento de la outbox}
\UASlideTag{NÚCLEO}{UAIngNavy}
\scriptsize
\begin{center}
\begin{tabular}{p{0.20\textwidth}p{0.64\textwidth}}
\toprule
Término & Definición \\
\midrule
Relay & Proceso que toma eventos pendientes y los publica en el broker. \\
Polling & El relay consulta periódicamente la tabla outbox. \\
CDC & \textit{Change Data Capture}: captura cambios confirmados en una base. \\
WAL / binlog & Registros internos desde los cuales una herramienta CDC puede leer cambios. \\
\bottomrule
\end{tabular}
\end{center}
\begin{alertblock}{Límite}
El relay puede perder el ACK y publicar otra vez. La outbox evita “DB sí / evento no”; no elimina duplicados downstream.
\end{alertblock}
\end{frame}

\begin{frame}{Debezium: un relay basado en CDC}
\UASlideTag{CASO REAL}{UAIngOrange}
\small
\begin{uainfo}
\textbf{Debezium} es una plataforma de conectores CDC. Puede leer cambios confirmados desde WAL o binlog y transformarlos en eventos publicados, frecuentemente en Kafka.
\end{uainfo}
\begin{center}
\begin{tikzpicture}[node distance=0.30cm]
\node[uaOrange,text width=1.65cm] (db) {DB Orders};
\node[uaSoft,text width=1.65cm,right=of db] (wal) {WAL / binlog};
\node[uaSoft,text width=1.75cm,right=of wal] (deb) {Debezium};
\node[uaOrange,text width=1.75cm,right=of deb] (top) {Topic};
\draw[uaOrangeArrow] (db)--(wal);
\draw[uaArrow] (wal)--(deb);
\draw[uaOrangeArrow] (deb)--(top);
\end{tikzpicture}
\end{center}
\begin{uakeybox}
Debezium es una implementación posible del relay; la idea arquitectónica es la outbox durable más publicación repetible.
\end{uakeybox}
\end{frame}

\begin{frame}{Flujo completo de 9F2A: outbox, broker e inbox}
\UASlideTag{NÚCLEO}{UAIngNavy}
\scriptsize
\begin{center}
\begin{tikzpicture}[node distance=0.22cm]
\node[uaOrange,text width=1.40cm] (api) {Checkout};
\node[uaSoft,text width=1.65cm,right=of api] (db) {DB\\pedido +\ outbox};
\node[uaSoft,text width=1.35cm,right=of db] (rel) {Relay};
\node[uaSoft,text width=1.40cm,right=of rel] (br) {Broker\\P2};
\node[uaSoft,text width=1.55cm,right=of br] (con) {Inventario};
\node[uaOrange,text width=1.55cm,right=of con] (eff) {DB\\inbox+stock};
\draw[uaOrangeArrow] (api)--(db);
\draw[uaArrow] (db)--(rel);
\draw[uaArrow] (rel)--(br);
\draw[uaArrow] (br)--(con);
\draw[uaOrangeArrow] (con)--(eff);
\end{tikzpicture}
\end{center}
\begin{itemize}
\item Outbox hace durable la intención de publicar.
\item Broker conserva y permite replay.
\item Inbox hace idempotente el efecto local.
\item El mismo \texttt{event\_id} conecta el recorrido extremo a extremo.
\end{itemize}
\end{frame}

\begin{frame}{Exactly-once: declare el objeto y la frontera}
\UASlideTag{NÚCLEO}{UAIngNavy}
\small
\begin{uakeybox}
“Exactly-once” no es una propiedad universal. Debe especificar qué estado, outputs y progreso se confirman juntos.
\end{uakeybox}
\begin{columns}[T]
\column{0.48\textwidth}
\begin{block}{Dentro de la frontera}
Input, estado local, output y progreso pueden participar en una transacción coordinada.
\end{block}
\column{0.48\textwidth}
\begin{alertblock}{Fuera de la frontera}
Email, pago o API externa requieren identidad, idempotencia, deduplicación o reconciliación propia.
\end{alertblock}
\end{columns}
\end{frame}

% =============================================================================
\UASectionSlide{Capítulo 5 --- operar un stream que nunca termina}
% =============================================================================

\begin{frame}{Stream y pipeline: Analítica lee una historia sin fin}
\UASlideTag{NÚCLEO}{UAIngNavy}
\small
\begin{uainfo}
\textbf{Stream:} secuencia potencialmente ilimitada de eventos.\\
\textbf{Pipeline:} cadena de etapas que filtran, transforman, agregan o enrutan esos eventos.
\end{uainfo}
\begin{center}
\begin{tikzpicture}[node distance=0.25cm]
\node[uaOrange,text width=1.55cm] (e) {Eventos};
\node[uaSoft,text width=1.55cm,right=of e] (f) {Filtrar};
\node[uaSoft,text width=1.55cm,right=of f] (g) {Agrupar};
\node[uaSoft,text width=1.55cm,right=of g] (a) {Agregar};
\node[uaOrange,text width=1.55cm,right=of a] (d) {Dashboard};
\draw[uaOrangeArrow] (e)--(f);
\draw[uaArrow] (f)--(g);
\draw[uaArrow] (g)--(a);
\draw[uaOrangeArrow] (a)--(d);
\end{tikzpicture}
\end{center}
\end{frame}

\begin{frame}{Stateless y stateful: qué debe recordar Analítica}
\UASlideTag{NÚCLEO}{UAIngNavy}
\small
\begin{columns}[T]
\column{0.48\textwidth}
\begin{block}{Stateless}
Transforma cada evento sin depender de eventos anteriores: normalizar moneda o proyectar campos.
\end{block}
\column{0.48\textwidth}
\begin{block}{Stateful}
Necesita memoria entre eventos: contar ventas por minuto, sumar ingresos o mantener ranking.
\end{block}
\end{columns}
\begin{uakeybox}
El estado puede vivir en RAM para velocidad, en SSD local para recuperación rápida y en un changelog remoto para reconstrucción si se pierde el nodo.
\end{uakeybox}
\end{frame}

\begin{frame}{Event time y processing time: el pedido llegó tarde}
\UASlideTag{NÚCLEO}{UAIngNavy}
\small
\begin{uainfo}
\textbf{Event time:} cuándo ocurrió el hecho en la fuente.\\
\textbf{Processing time:} cuándo lo procesa la etapa actual.
\end{uainfo}
\begin{center}
\begin{tikzpicture}[x=1.0cm,y=0.75cm]
\draw[->,thick,UAIngNavy] (0,0)--(9.1,0) node[right,uaTiny]{tiempo};
\node[uaEventOrange] at (1.4,0.55) {9F2A ocurrió 12:00};
\node[uaEvent] at (5.0,-0.55) {broker 12:02};
\node[uaEventOrange] at (7.3,0.55) {procesado 12:03};
\end{tikzpicture}
\end{center}
\begin{alertblock}{Consecuencia}
Agrupar por tiempo de procesamiento puede colocar una venta de las 12:00 en la ventana de las 12:03.
\end{alertblock}
\end{frame}

\begin{frame}{Watermark: cuándo Analítica decide cerrar una ventana}
\UASlideTag{EXTENSIÓN}{UAIngNavy}
\small
\begin{uainfo}
Un \textbf{watermark} es una estimación de progreso en event time: “no espero normalmente eventos anteriores a este instante”.
\end{uainfo}
\begin{itemize}
\item No garantiza que jamás llegue un evento más antiguo.
\item Permite cerrar ventanas sin esperar para siempre.
\item Un evento posterior al watermark es late data.
\item La política puede descartar, corregir, reabrir o enviar a una salida lateral.
\end{itemize}
\begin{uakeybox}
La política depende del costo del error: un dashboard puede corregirse; una decisión antifraude quizá espere más.
\end{uakeybox}
\end{frame}

\begin{frame}{Transferencia: ranking de Asteria Online}
\UASlideTag{ACTIVIDAD}{UAIngOrange}
\small
Diseñe el ranking de combate:
\begin{enumerate}
\item evento y \texttt{event\_id};
\item key y orden que necesita;
\item número inicial de particiones;
\item estado que conserva el consumidor;
\item tratamiento de eventos tardíos;
\item SLO de atraso;
\item prueba de crash y replay.
\end{enumerate}
\begin{alertblock}{Pregunta de transferencia}
¿Qué cambiaría si el ranking entrega premios irreversibles al cerrar la temporada?
\end{alertblock}
\end{frame}

\begin{frame}{La historia completa de 9F2A y 1C7B}
\UASlideTag{NÚCLEO}{UAIngNavy}
\scriptsize
\begin{enumerate}
\item Checkout confirma cada pedido y guarda outbox.
\item Relay publica eventos con \texttt{key=order\_id}.
\item 9F2A llega a P2 y 1C7B a P0.
\item Cada partición preserva la historia de sus pedidos.
\item Inventario reparte P0, P1 y P2 entre instancias.
\item Si una instancia cae, otra retoma desde committed offset.
\item Inbox evita repetir stock ante replay.
\item Analítica procesa el mismo log con su propio progreso y política temporal.
\end{enumerate}
\begin{uakeybox}
Los conceptos ya no forman una lista: describen etapas y fallas del mismo recorrido extremo a extremo.
\end{uakeybox}
\end{frame}

\begin{frame}{Método reusable para diseñar un pipeline}
\UASlideTag{NÚCLEO}{UAIngNavy}
\small
Para cada flujo, responda en este orden:
\begin{enumerate}
\item ¿Cuál es el hecho y cuál su identidad estable?
\item ¿Qué eventos requieren orden relativo?
\item ¿Qué key agrupa esa causalidad sin concentrar carga?
\item ¿Qué particiones y consumidores necesita el throughput?
\item ¿Dónde se guarda el progreso recuperable?
\item ¿Qué efecto debe ser idempotente?
\item ¿Existe dual write? ¿Outbox/inbox?
\item ¿Qué métricas detectan skew, lag y rebalances?
\end{enumerate}
\end{frame}

\begin{frame}{Volvemos al Diagnóstico inicial}
\UASlideTag{ACTIVIDAD}{UAIngOrange}
\small
Responda ahora con vocabulario técnico:
\begin{enumerate}
\item ¿Por qué \texttt{order\_id} preserva la causalidad de cada pedido?
\item ¿Por qué \texttt{campus\_id} puede producir una hot partition?
\item ¿Qué reaparece si el efecto fue durable pero el committed offset no?
\item ¿Qué resuelve outbox y qué deja abierto?
\item ¿Por qué agregar consumidores no divide una partición?
\end{enumerate}
\end{frame}

\begin{frame}{Síntesis: cinco ideas que deben quedar}
\UASlideTag{NÚCLEO}{UAIngNavy}
\small
\begin{enumerate}
\item La key define el ámbito de orden y condiciona la distribución de carga.
\item Una partición ofrece orden local y es unidad de almacenamiento y paralelismo.
\item Committed offset registra progreso; no confirma por sí solo el efecto de negocio.
\item Replay es normal; idempotencia e inbox vuelven seguro el efecto.
\item Outbox conecta una transacción local con publicación repetible, no con exactly-once universal.
\end{enumerate}
\end{frame}

\begin{frame}{Exit ticket}
\UASlideTag{ACTIVIDAD}{UAIngOrange}
\small
En una hoja, responda:
\begin{enumerate}
\item Defina una buena key en una frase y aplíquela a 9F2A.
\item Dibuje dónde quedan 9F2A y 1C7B en tres particiones.
\item Explique qué ocurre si Inventario cae después del efecto y antes del committed offset.
\item Nombre una métrica que distinga hot partition de atraso general.
\end{enumerate}
\end{frame}

\begin{frame}{Fuentes y material de respaldo}
\UASlideTag{RESPALDO}{UAIngNavy}
\scriptsize
\begin{itemize}
\item Apache Kafka Documentation: topics, partitions, producer, consumer groups y delivery semantics.
\item Kreps et al.: \emph{Kafka: a Distributed Messaging System for Log Processing}.
\item Kleppmann: \emph{Designing Data-Intensive Applications}, capítulos sobre logs, particionado y stream processing.
\item Debezium Documentation: transactional outbox y CDC.
\item Apache Flink Documentation: event time, watermarks y stateful processing.
\end{itemize}
\begin{uakeybox}
Lea siempre la documentación buscando cinco preguntas: orden, progreso, falla, efecto y costo operacional.
\end{uakeybox}
\end{frame}

\end{document}
